\chapter{Introduction} \label{chap:introduction}

\section{Overview}

Imagine you're a nurse on a night shift, trying to decide whether a new treatment protocol for influenza is worth recommending to a patient. Or you're a college student picking a major, trying to understand what problems your field is actually working on. Or maybe you're simply a parent, watching the news about a new coronavirus, and you want to know what it means for your family. Each of these readers is making a decision that turns on what scientific evidence actually says, and how well they understand it shapes how well they decide. The link is most clearly documented in medicine, where a person's ability to find, understand, and act on health information is consistently associated with concrete outcomes. Lower health literacy is linked to higher mortality, more preventable hospitalizations, and poorer treatment adherence~\cite{berkman2011lowhealthliteracy, paasche-orlow2007causal}. The same logic extends beyond the clinic. A student who can't parse a field's literature can't choose it. A voter who can't follow an evidence debate is left to defer to whoever sounds most confident.
 
These individual stakes aggregate into a societal one. Public understanding of science underwrites trust in scientific institutions and the willingness of policymakers and citizens to act on evidence, especially during crises~\cite{cologna2025trust, intemann2023scicomm}. That trust has been eroding. Pew Research Center surveys show that the share of U.S. adults expressing at least a fair amount of confidence in scientists to act in the public's best interest fell from 87\% in April 2020 to 73\% by 2023, with the strongest level of trust nearly halving over the same period~\cite{pew2023trust, pew2024trust}. The consequences are not merely attitudinal. The public doesn't just benefit from science, they help pay for it. Around 40\% of basic research in the U.S. is government funded~\cite{Pece_Anderson_2024}. Recent federal budget actions have begun to reflect this erosion. In 2025 alone, more than 3,800 active grants at NIH and NSF, totaling roughly \$3 billion in unspent funds, were frozen or terminated~\cite{sciencenews2025cuts, Agarwal_2025}. Whatever the proximate political causes, the broader pattern is a public that is less able to see what science produces becomes less willing to keep paying for it. Good scientific communication is essential for building public trust.
 
But open access is not the same as accessible. A paper freed from a paywall is still locked behind years of specialized training, and even the lay summaries written to bridge that gap often miss their mark. Systematic reviews of plain language summaries in health journals find that most are still written well above recommended reading levels and rely on jargon despite explicit guidance to avoid it~\cite{Gainey2024ArePL}. More fundamentally, every one of these efforts, produces a single artifact aimed at a single imagined reader. The nurse, the student, and the concerned citizen come to the same paper with different questions, different background knowledge, and different stakes. A summary that works for a clinician will overwhelm a high schooler; a summary that works for a high schooler will be unusable to a researcher. A summary cannot serve them all.
 
Large language models open up a different approach. Their ability to follow flexible, natural-language instructions and to adapt style and content on the fly means that personalization no longer requires building a separate system for every user type. Recent work has used LLMs to generate plain language summaries of biomedical research~\cite{Guo2020AutomatedLL,luo2022readability}. However, this work largely inherits the same generic-reader assumption from earlier summarization systems: a fixed reading level, a fixed level of detail, a fixed notion of what counts as ``important.'' These assumptions make systems easier to build and evaluate, but they also limit how useful those systems can be. The question is not how to produce \emph{the} accessible version of a paper, but rather what does this particular reader need, and how do we build and evaluate systems that deliver it?
 
\textbf{The core hypothesis of this thesis is that we can achieve better summaries by centering the needs of the end user in our generation and evaluation methodologies.}
 
\Cref{part1} argues that summaries are better when the goals and interests of the reader are modeled explicitly. We focus on 2 axes of controllability: structure and readability. \Cref{chap:chap-3} introduces a method for user-defined summary structure. Many documents, such as blogs, posters, or slides, can be considered structured summaries of an underlying scientific paper. The generative framework works in a 2 step process. First we induce structure into the knowledge represented in a scientific paper, then use that representation to generate the final summary with a user-indicated structure. We then extend the controllability along the axis of readability in~\Cref{chap:chap-4}. We introduce a light-weight, training-free method for fine-grained controllable generation, letting users specify desired attributes along a continuous spectrum of readability rather than assuming readability to be a binary choice between lay and technical. We show that control vectors, along with a requirements-based framework for data selection, generate high-quality, controllable summaries in a low-compute setting. We focus the methods and experiments in~\Cref{part1} on smaller, lower-cost models to increase the accessibility of our methodology.
 
\Cref{part2} argues that in generating personalized summaries, the evaluation must also account for the user. \Cref{chap:chap-5} shows that traditional readability metrics correlate poorly with human judgments of readability. We find that the most popular metrics in plain language summarization literature suffer from both definitional inconsistency and measurement error. In other words, they are not good measures of readability for plain language summarization tasks. We show that LLMs capture deeper attributes of readability, such as required background knowledge. We recommend that future work in plain language summarization focus their evaluation on LLM-as-judges, higher performing traditional metrics, and qualitative analysis. Through the goal of evaluating personalized summaries, we found that there is little gold data for personalized generation, and that collecting this data is prohibitively costly given the increase in requirements for implementation. We address this problem in~\Cref{chap:chap-6}, in which we leverage naturally occurring data and cheaply acquired machine-generated summaries to improve the performance of LLMs-as-judges. Finally, in~\Cref{chap:chap-7}, we use persona-based data to measure whether a summary actually meets a particular user's informational needs. We show that the majority of traditional metrics for summarization are not sufficiently robust to changes in informational content, and that the performance of current LLMs does not sufficiently measure personalized information content.
 
Together, these chapters make the case that the next generation of summarization systems should be designed and evaluated around the people who will use them, not around an abstract average reader.

\section{Publications}

Much of the work in this thesis is either published or under-review for publication. Below is a list of the publications completed during my graduate program, and a brief summary of each.\footnote{The irony of writing general summaries in a thesis about personalized summaries is not lost on me.}

\subsection{Summarization Research}
\begin{itemize}
    \item We improve generation of structured summaries by building an intermediate knowledge representation and introduce a evaluation framework that accounts for the intended structure of the summary~\cite[][\Cref{chap:chap-3}]{cachola-etal-2024-knowledge}.
    \item We use citation networks to generate explanations on the relationships between scientific publications~\cite{luu-etal-2021-explaining}.
    \item We leverage the latest research in AI and HCI to build an interactive reading interface for scholarly works~\cite{semantic-reader-project}.
    \item We show that the most popular metrics used to evaluate readability are poorly correlated with human judgments and we show that LLMs measure deeper attributes of readability, such as explanations of technical concepts~\cite[][\Cref{chap:chap-5}]{cachola-etal-2025-evaluating}. 
    % \item  We propose a light weight control vector method for fine-grained readability control in scientific summarization along with a requirements-based framework for data selection~[\todocite,\Cref{chap:chap-4}].
    % \item We show that LLMs can be better judges of summaries by leveraging related but non-exact references~[\todocite,\Cref{chap:chap-6}].

\end{itemize}
\subsection{Additional Contributions}
In addition to my core work on summarization research, listed above, I have had the privilege to contribute to work in a variety of domains during my graduate career. These contributions are listed and briefly summarized below.
\begin{itemize}
    \item We build a platform that converts PDFs to HTML, with the focus on making science more accessible to blind or low-vision readers~\cite{wang-cachola-scially}. {\small \badge[blue]{Best Artifact Award}}
    \item We conduct an analysis of the effects of shot selection for in-context learning on the fairness of LLMs, and show that overall performance is not indicative of their performance~\cite{aguirre-etal-2024-selecting}.
    \item We use knowledge distillation to generate faithful and plausible explanations of medical code predictions~\cite{wood-doughty-etal-2022-model}.
    \item We introduce the task of novel dog-whistle discovery and present a strong baseline system that leverages the strengths of vector databases and Large Language Models~\cite{sasse-etal-2025-making}.
\end{itemize}

\section{Non-technical abstracts}

I am a teacher at heart. While a PhD dissertation is often (necessarily) technical in nature, my hope is that anyone who reads this thesis, regardless of background, can walk away with having learned something new. It would feel disingenuous to write a body of work about designing systems with the goal of making science accessible, and that work be inaccessible. To this end, I will include a Non-technical abstract at the beginning of each chapter. This abstract will summarize the main motivations, contributions, and findings, so that anyone of any background can understand the main take-aways from each chapter.